\documentclass[11pt,a4paper]{article}
\usepackage{amsmath,amssymb,amsthm}
\usepackage{mathtools}
\usepackage{bm}
\usepackage{booktabs}
\usepackage{array}
\usepackage{enumitem}
\usepackage{geometry}
\usepackage{graphicx}
\usepackage{xcolor}
\usepackage{tcolorbox}
\usepackage[pdfencoding=auto,psdextra]{hyperref}

\geometry{margin=2.5cm}

\newcommand{\E}{\mathbb{E}}
\newcommand{\R}{\mathbb{R}}
\newcommand{\norm}[1]{\left\lVert #1 \right\rVert}
\newcommand{\Fc}{\mathcal{F}_C}
\newcommand{\uvec}{\bm{u}}
\newcommand{\xvec}{\bm{x}}

\newtheorem{proposition}{Proposition}
\newtheorem{theorem}{Theorem}
\newtheorem{corollary}{Corollary}
\newtheorem{definition}{Definition}
\newtheorem*{remark*}{Remark}

\title{\textbf{Mixture-of-Flow velocity decomposition:}\\[0.3em]
\large capacity ceiling vs.\ optimization gap, the posterior-responsibility form of a
multi-domain teacher velocity, and when a MoF can be split into valid sub-flows}
\author{Flow Factory}
\date{}

\begin{document}
\maketitle

\begin{abstract}
We formalize the Mixture-of-Flow (MoF) student for cross-model on-policy distillation (XOPD): a set
of full flow-matching experts whose output velocities are blended by a router. We prove that the
marginal velocity of a \emph{mixture-of-domains} teacher is \emph{exactly} a posterior-responsibility
weighted convex combination of the per-domain velocities (Thm.~\ref{thm:posterior}), which pins down
the ``correct'' functional form of a MoF student and explains why a sum-to-one (convex) gate makes
every expert an individually valid flow. We then resolve an apparent paradox
(\enquote{$\text{big}\!\to\!\text{MoF}\!\to\!\text{single}$ beats $\text{big}\!\to\!\text{single}$})
by separating the \emph{representational ceiling} (capacity-bound, path-independent) from the
\emph{optimization gap} (path-dependent). Finally we state when a mixture with fixed \emph{activated}
capacity can exceed a single model of the same activated capacity, and contrast MoF with classic
(FFN-level) MoE. Companion progress note:
\texttt{docs/xopd/progress/2026-07-21-mof-velocity-decomposition.md}.
\end{abstract}

\section{Setup and notation}
Flow matching learns a time-dependent velocity field $\uvec_t:\R^d\to\R^d$, $t\in[0,1]$, whose ODE
$\dot\xvec_t=\uvec_t(\xvec_t)$ transports a base density $p_0$ (noise) to the data density $p_1$. For a
probability path $(p_t)$ induced by a coupling with per-sample conditional velocity
$\uvec_t(\xvec\mid z)$ (where $z$ is the data-side latent/condition), the \emph{marginal} velocity is a
conditional expectation
\begin{equation}\label{eq:marginal}
  \uvec_t(\xvec)=\E\!\left[\dot\xvec_t \,\middle|\, \xvec_t=\xvec\right]
             =\int \uvec_t(\xvec\mid z)\, p_t(z\mid \xvec)\,\mathrm{d}z .
\end{equation}
The map ``target $\mapsto$ marginal velocity'' is \textbf{linear} (it is an expectation); this is the
single structural fact behind everything below.

\begin{definition}[Valid flow]
A field $\uvec_t$ is a \emph{valid flow} for a path $(p_t)$ if it satisfies the continuity equation
$\partial_t p_t + \nabla\!\cdot(p_t\,\uvec_t)=0$ with boundary $p_0=$ noise, $p_1=$ data; equivalently
it is the marginal velocity \eqref{eq:marginal} of some coupling reaching $p_1$.
\end{definition}

\section{The teacher velocity of a multi-domain mixture is a posterior-weighted convex blend}

Let the data be a mixture over $N$ domains (e.g.\ \texttt{geneval}, \texttt{ocr}),
$q=\sum_{e=1}^N \lambda_e\, q_e$, $\lambda_e>0$, $\sum_e\lambda_e=1$, each domain $e$ inducing its own
path $p_t^{(e)}$ and valid flow $\uvec_t^{(e)}$.

\begin{theorem}[Posterior-responsibility decomposition]\label{thm:posterior}
The marginal velocity of the mixture teacher is
\begin{equation}\label{eq:posterior}
  \uvec_t(\xvec)=\sum_{e=1}^N \pi_e(\xvec,t)\,\uvec_t^{(e)}(\xvec),
  \qquad
  \pi_e(\xvec,t)=\frac{\lambda_e\, p_t^{(e)}(\xvec)}{\sum_{e'}\lambda_{e'}\, p_t^{(e')}(\xvec)},
\end{equation}
where $\pi_e(\xvec,t)\ge 0$ and $\sum_e \pi_e(\xvec,t)=1$ is the Bayes posterior that $\xvec_t$ was
produced by domain $e$.
\end{theorem}

\begin{proof}
The mixture path density is $p_t(\xvec)=\sum_e \lambda_e p_t^{(e)}(\xvec)$. Introduce the domain
indicator $E$ with $P(E=e)=\lambda_e$. By the tower rule applied to \eqref{eq:marginal},
\[
  \uvec_t(\xvec)=\E[\dot\xvec_t\mid \xvec_t=\xvec]
  =\sum_e P(E=e\mid \xvec_t=\xvec)\,\E[\dot\xvec_t\mid \xvec_t=\xvec, E=e].
\]
The inner conditional expectation is by definition the per-domain marginal velocity
$\uvec_t^{(e)}(\xvec)$, and Bayes' rule gives
$P(E=e\mid \xvec_t=\xvec)=\lambda_e p_t^{(e)}(\xvec)/\sum_{e'}\lambda_{e'}p_t^{(e')}(\xvec)=\pi_e(\xvec,t)$.
\end{proof}

\begin{corollary}[Correct MoF form]\label{cor:form}
A per-sample MoF with $N$ experts, softmax (sum-to-one) gate $w_e(\xvec,t)$ and blend
$\hat\uvec_t=\sum_e w_e\,\hat\uvec_t^{(e)}$ can represent the multi-domain teacher \emph{exactly} in the
infinite-capacity limit, by matching $\hat\uvec^{(e)}\!\to\!\uvec^{(e)}$ (per-domain flow) and
$w_e\!\to\!\pi_e$ (posterior). Hence the convex constraint $\sum_e w_e=1$ is not a heuristic: it is the
structural form of \eqref{eq:posterior}.
\end{corollary}

\begin{remark*}
Eq.~\eqref{eq:posterior} is a \emph{convex} combination with input-dependent weights. The intuition
``a velocity field is a sum of simpler velocity fields'' is precise only in this posterior-convex
sense, \textbf{not} as a free additive sum (\S\ref{sec:validity}).
\end{remark*}

\section{Validity of the split: convex vs.\ free gates}\label{sec:validity}

\begin{proposition}[When each expert is a valid flow]\label{prop:validity}
Let $\hat\uvec_t=\sum_e w_e(\xvec,t)\,\hat\uvec_t^{(e)}$.
\begin{enumerate}[leftmargin=1.4em]
  \item \textbf{Convex, $\sum_e w_e=1$, $w_e\ge0$ (softmax / hard routing):} if the experts match the
  per-domain flows of Thm.~\ref{thm:posterior}, each $\hat\uvec^{(e)}$ is individually a valid flow
  and the blend is the valid mixture flow. Under \emph{hard} routing ($w\in\{0,1\}$, disjoint domains)
  the blend is piecewise-single-expert.
  \item \textbf{Free / additive (e.g.\ independent sigmoid gates, $\sum_e w_e\neq1$):} the blend may
  fit a strictly larger function class, but an individual summand $\hat\uvec^{(e)}$ need not be a valid
  flow (it can be a partial/curl-like component that alone does not transport $p_0\!\to\!p_1$), and the
  blend need not correspond to any single clean probability path.
\end{enumerate}
\end{proposition}

Thus \textbf{sum-to-one trades capacity for decomposability}: it is the exact structure that keeps
``each expert = a valid flow'' and makes ``split the MoF back into standalone flows'' meaningful; free
gates buy raw fit at the cost of that interpretation. This is the crux of the ``keep vs.\ drop the
convex constraint'' design choice.

\section{The apparent paradox and its resolution}

The informal argument is: (i) $\text{big}\!\to\!\text{single}$ is lossy; (ii) $\text{big}\!\to\!\text{MoF}$
is lossy but less; (iii) $\text{MoF}\!\to\!\text{single}$ is \emph{lossless}; therefore
$\text{big}\!\to\!\text{MoF}\!\to\!\text{single}$ beats $\text{big}\!\to\!\text{single}$, contradicting a
capacity-determined ceiling and the expectation that staging loses performance.

\paragraph{Decompose achieved loss.}
For a hypothesis class $\mathcal F$ and algorithm $\mathcal A$, write the achieved loss as
\begin{equation}\label{eq:decomp}
  \underbrace{\mathrm{Loss}(\mathcal A,\mathcal F)}_{\text{achieved}}
  = \underbrace{R(\mathcal F)}_{\text{representational ceiling}} + \underbrace{G(\mathcal A,\mathcal F)}_{\text{optimization gap}},
  \qquad R(\mathcal F)=\inf_{f\in\mathcal F}\mathrm{Loss}(f).
\end{equation}
Let $\Fc$ be the single-base class (capacity $C$) and $\mathcal F_{\text{MoF}}\supseteq\Fc$ the MoF class.

\begin{proposition}[The ``lossless'' leg is regime-dependent]\label{prop:regime}
$\text{MoF}\!\to\!\text{single}$ is (near) lossless \emph{iff} the MoF output function lies in $\Fc$.
\begin{itemize}[leftmargin=1.4em]
  \item \textbf{Disjoint-domain / hard routing:} $\pi_e\in\{0,1\}$, so per input only one expert acts;
  the MoF output is piecewise-single-expert and lies in $\Fc$ whenever $C$ suffices for the
  \emph{union} of per-domain functions. Then $\text{MoF}\!\to\!\text{single}$ is near-lossless.
  \item \textbf{Overlapping domains / soft routing:} $\pi_e\in(0,1)$ forces a genuine blend of two
  velocities at the \emph{same} input; the output generally lies \emph{outside} $\Fc$, so
  $\text{MoF}\!\to\!\text{single}$ is \textbf{lossy}.
\end{itemize}
\end{proposition}

\begin{theorem}[No capacity is created by staging]\label{thm:ceiling}
The final student is a single base in $\Fc$, hence its achievable loss is bounded below by $R(\Fc)$
regardless of the path. In particular
$\mathrm{Loss}\big(\text{big}\!\to\!\text{MoF}\!\to\!\text{single}\big)\ge R(\Fc)$, the same
representational floor as $\text{big}\!\to\!\text{single}$.
\end{theorem}

\begin{corollary}[Resolution]\label{cor:resolution}
The three premises never hold simultaneously in a way that beats $R(\Fc)$:
\begin{itemize}[leftmargin=1.4em]
  \item In the disjoint regime, leg (iii) is $\approx$lossless, but then leg (ii)'s advantage over
  $\text{big}\!\to\!\text{single}$ is \emph{not} a higher ceiling (a single base of the same per-expert
  capacity has the same floor on the union); it is an \emph{optimization} advantage (no cross-domain
  interference during fitting), i.e.\ a reduction of $G$, not of $R$.
  \item In the overlapping regime, leg (ii) genuinely lowers $R$ (MoF has more capacity), but leg
  (iii) is then lossy: compressing a $\mathcal F_{\text{MoF}}$ function back into $\Fc$ pays exactly the
  capacity difference.
\end{itemize}
Either way, staging cannot beat the $\Fc$ ceiling. What staging \emph{can} do is shrink the
optimization gap $G(\mathcal A,\Fc)$: the well-optimized MoF is a smoother, easier distillation target
(dark knowledge / born-again / progressive distillation), letting the single base approach its
\emph{own} ceiling $R(\Fc)$ more closely than direct training does. Thus
$\text{big}\!\to\!\text{MoF}\!\to\!\text{single}$ can beat \emph{direct} $\text{big}\!\to\!\text{single}$
\textbf{without contradiction}: it recovers optimization gap, not representational capacity.
\end{corollary}

\section{Fixed activated capacity: when does a mixture beat a single base?}

Consider top-$k$-of-$N$ routing: \emph{activated} parameters per input are $\approx kC/1$ (here we take
$k{=}1$ so activated $=C$), while \emph{total} parameters are $NC$.

\begin{proposition}[Partitioning gain]\label{prop:partition}
With activated capacity fixed at $C$:
\begin{itemize}[leftmargin=1.4em]
  \item If the teacher velocity is \emph{routable} (input space splits into regions on which the
  target is simpler, e.g.\ multi-domain or a nontrivial posterior $\pi(\xvec,t)$), then $N$ experts each
  dedicate the full $C$ to one region, whereas a single $C$ must \emph{share} $C$ across all regions
  (cross-region interference). The mixture's effective ceiling is the \emph{region-wise} best, which
  dominates the shared-$C$ ceiling: $R(\mathcal F_{\text{MoF},\,k{=}1})\le R(\Fc)$, often strictly.
  \item If the teacher is \emph{monolithic} (per-input intrinsic complexity $>C$ everywhere), each
  input still sees only $C$ activated parameters, so the mixture cannot exceed the single-$C$ ceiling.
\end{itemize}
\end{proposition}

\noindent Hence a mixture of identical small bases can match a large teacher at constant activated
compute \emph{only} to the extent the teacher's velocity is routable/decomposable; this is precisely
the structure Thm.~\ref{thm:posterior} exhibits for multi-domain data.

\section{Contrast with classic (FFN-level) MoE}

\begin{center}
\begin{tabular}{@{}lll@{}}
\toprule
 & \textbf{Classic MoE} & \textbf{MoF (this work)} \\
\midrule
Expert unit & FFN sub-module (not a flow) & full transformer $=$ a valid flow \\
Routing & per-token, per-layer (fine) & per-sample, output-level (coarse) \\
Mixing expressivity & high (re-routes each layer) & low (linear blend of final velocities) \\
Split into standalone flows & no & yes (convex case, Prop.~\ref{prop:validity}) \\
Param efficiency & high & low (coarse experts) \\
Matches Thm.~\ref{thm:posterior} form & implicitly, entangled & explicitly ($w_e\!\approx\!\pi_e$) \\
\bottomrule
\end{tabular}
\end{center}

\noindent Classic MoE is the stronger tool for \emph{fitting} a teacher at fixed activated compute
(Prop.~\ref{prop:partition}); MoF is the right object for the \emph{scientific} question of decomposing
a teacher into individually valid, recombinable sub-flows (Cor.~\ref{cor:form},
Prop.~\ref{prop:validity}). They optimize different objectives: fit efficiency vs.\ decomposability.

\section{Testable predictions}
\begin{enumerate}[leftmargin=1.4em]
  \item \textbf{Domain overlap is the master switch.} The more disjoint the domains (posterior
  $\pi\!\to\!\{0,1\}$, measurable via the router-specialization visualization), the more
  $\text{big}\!\to\!\text{MoF}$'s gain is \emph{optimization} and the closer
  $\text{MoF}\!\to\!\text{single}$ is to lossless; more overlap $\Rightarrow$ real ceiling gain but lossy
  merge.
  \item \textbf{Staging beats direct only up to $R(\Fc)$.} $\text{big}\!\to\!\text{MoF}\!\to\!\text{single}$
  should modestly beat direct $\text{big}\!\to\!\text{single}$, by an amount $\approx$ the direct
  optimization gap; a \emph{large} win indicates a large direct optimization gap (itself worth study),
  not broken capacity.
  \item \textbf{Convex vs.\ free gate.} Free (sigmoid) gates should reach $\le$ training MSE($v$) of
  convex (softmax) gates, but the \emph{standalone} sampling quality of an individually extracted
  expert should be higher for convex (Prop.~\ref{prop:validity}). This separates ``fit'' from
  ``valid decomposition''.
  \item \textbf{Fair baseline.} A $k{=}1$ MoF (activated $C$) must be compared against a single base of
  capacity $C$ (equal activated compute), and against a classic MoE of equal activated compute, not
  against a $2C$ dense model.
\end{enumerate}

\end{document}
